% Part: normal-modal-logic
% Chapter: axioms-systems
% Section: derived-rules

\documentclass[../../../include/open-logic-section]{subfiles}

\begin{document}

\olfileid[hi]{nml}{prf}{der}

\olsection{व्युत्पन्न नियम}

!!{derivation}याँ खोजना और लिखना स्पष्टतः कठिन, बोझिल और पुनरावृत्तिपूर्ण
है। उदाहरणार्थ, हम बहुत बार $!A \lif !B$ से
$\Box !A \lif \Box !B$ तक जाना, अर्थात् नियम~\RK{} लागू करना चाहते हैं।
इसके लिए पहले \Nec{} का अनुप्रयोग, फिर~\Ax{K} का उचित दृष्टांत दर्ज करना,
फिर~\MP{} लागू करना पड़ता है। $!A \lif !B$ और $!B \lif !C$ से
$!A \lif !C$ तक जाने के लिए यह (लंबा) सर्वसत्यतात्मक दृष्टांत दर्ज करना
पड़ता है:
\[
(!A \lif !B) \lif ((!B \lif !C) \lif (!A \lif !C))
\]
और \MP{} दो बार लागू करना पड़ता है। हम प्रायः किसी उप-!!{formula} को ऐसे
सूत्र से बदलना चाहते हैं जिसे हम उसके तुल्य जानते हैं, जैसे
\iftag{prvDiamond}{$\Diamond !A$ को $\lnot\Box\lnot !A$ से}{$\Box !A$ को
$\lnot\Diamond\lnot !A$ से}, या $\lnot\lnot !A$ को $!A$ से। इसलिए वास्तविक
!!{derivation} लिखने के स्थान पर केवल यह दर्ज करना अधिक सुविधाजनक है कि
मध्यवर्ती चरण क्यों !!{derivable} हैं। इस उद्देश्य से !!{derivability} के
कुछ तथ्य एकत्र करें।

\begin{prop}
  यदि $\Log{K} \Proves !A_1$, \dots, $\Log{K} \Proves !A_n$, और $!B$,
  $!A_1$, \dots, $!A_n$ से प्रतिज्ञप्तिक तर्कशास्त्र द्वारा निकलता है, तो
  $\Log{K} \Proves !B$।
\end{prop}

\begin{proof}
  यदि $!B$, $!A_1$, \dots,~$!A_n$ से प्रतिज्ञप्तिक तर्कशास्त्र द्वारा
  निकलता है, तो
  \[
  !A_1 \lif (!A_2 \lif \cdots (!A_n \lif !B)\dots)
  \]
  एक सर्वसत्यतात्मक दृष्टांत है। \MP{} को $n$ बार लागू करने से~$!B$ की
  !!{derivation} मिलती है।
\end{proof}

इस प्रतिज्ञप्ति के उपयोग को हम~\PL{} से दर्शाएँगे।

\begin{prop}
  यदि $\Log{K} \Proves !A_1 \lif (!A_2 \lif \cdots (!A_{n-1} \lif
  !A_n)\dots)$, तो $\Log{K} \Proves \Box !A_1 \lif (\Box !A_2 \lif
  \cdots (\Box !A_{n-1} \lif \Box !A_n)\dots)$।
\end{prop}

\begin{proof}
  $n$ पर आगमन से, ठीक \olref[nor]{prop:rk} के प्रमाण की तरह।
\end{proof}

इस प्रतिज्ञप्ति के उपयोग को हम~\RK{} से दर्शाएँगे। देखें कि ये परिणाम
!!{derivability} के परिणाम अधिक आसानी से स्थापित करने में कैसे सहायक हैं।

\begin{prop}
  $\Log{K} \Proves (\Box!A \land \Box!B) \lif \Box (!A \land !B)$
\end{prop}

\begin{proof}
  \begin{derivation}
    1. & $\Log{K} \Proves !A \lif (!B \lif (!A \land !B))$ & \Taut \\
    2. & $\Log{K} \Proves \Box!A \lif (\Box !B \lif \Box(!A \land !B)))$ & \RK, 1\\
    3. & $\Log{K} \Proves (\Box!A \land \Box!B) \lif \Box(!A \land !B)$ & \PL, 2
  \end{derivation}
\end{proof}

\begin{prop}\ollabel{prop:rewriting}
  यदि $\Log{K} \Proves !A \liff !B$ और $\Log{K} \Proves
  \Subst{!C}{!A}{q}$, तो $\Log{K} \Proves \Subst{!C}{B}{q}$
\end{prop}

\begin{proof}
  अभ्यास।
\end{proof}

\begin{prob}
  $!C$ की जटिलता पर आगमन से यह सिद्ध करके
  \olref[nml][prf][der]{prop:rewriting} सिद्ध कीजिए कि यदि
  $\Log{K} \Proves !A \liff !B$, तो
  $\Log{K} \Proves \Subst{!C}{!A}{q} \liff \Subst{!C}{!B}{q}$।
\end{prob}

यह प्रतिज्ञप्ति विशेष रूप से तब उपयोगी है जब हम $\Diamond$ को $\Box$ में
(या इसके विपरीत) बदलना, या किसी !!{formula} के भीतर द्विनिषेध हटाना चाहते
हैं। आगे जब भी हम !!{formula}~$!C(!A)$ को~$!C(!B)$ के रूप में पुनर्लिखेंगे,
\olref{prop:rewriting} के अनुप्रयोग को ``$!B$ के स्थान पर $!A$'' से चिह्नित
करेंगे। दूसरे शब्दों में, ``$!B$ के स्थान पर $!A$'' इसका संक्षिप्त रूप है:
  \begin{derivation}
    & $\Proves !C(!A)$ \\
    & $\Proves !A \liff !B$\\
    & $\Proves !C(!B)$ & \olref{prop:rewriting} से
  \end{derivation}
उदाहरणार्थ:

\begin{prop}
  $\Log{K} \Proves \lnot\Box p \lif \Diamond \lnot p$
\end{prop}

\begin{proof}
  \iftag{prvDiamond}{% Diamond is primitive
  \begin{derivation}
    1. & $\Log{K} \Proves \Diamond \lnot p \liff \lnot\Box\lnot\lnot p$ &
    \Dual\\
    2. & $\Log{K} \Proves \lnot \Box \lnot\lnot p \lif \Diamond \lnot p$ & \PL, 1\\
    3. & $\Log{K} \Proves \lnot\Box p \lif \Diamond \lnot p$ & $\lnot\lnot p$ के स्थान पर $p$\\
  \end{derivation}
  }{% Diamond is defined
  \begin{derivation}
    1. & $\Log{K} \Proves \lnot \Box \lnot\lnot p \lif \Diamond \lnot p$ & \Taut\\
    2. & $\Log{K} \Proves \lnot\Box p \lif \Diamond \lnot p$ & $\lnot\lnot p$ के स्थान पर $p$\\
  \end{derivation}
  पंक्ति~$1$ का सूत्र $\lnot q \lif \lnot q$ का वह दृष्टांत है जिसमें
  $q$ के स्थान पर $\Box\lnot\lnot p$ रखा गया है। $\Diamond\lnot p$ और
  $\lnot\Box\lnot\lnot p$ सर्वसम हैं, क्योंकि $\Diamond$ को
  $\lnot\Box\lnot$ के रूप में परिभाषित किया गया है।}
\end{proof}

ऊपर की !!{derivation} में अंतिम चरण ``$\lnot\lnot p$ के स्थान पर $p$''
इसका संक्षिप्त रूप है:
  \begin{derivation}
    & $\Log{K} \Proves \lnot \Box \lnot\lnot p \lif \Diamond \lnot
    p$\\
    & $\Log{K} \Proves \lnot\lnot p \liff p$ & \Taut\\
    & $\Log{K} \Proves \lnot\Box p \lif \Diamond \lnot p$ & \olref{prop:rewriting} से
  \end{derivation}
\olref{prop:rewriting} के $!C(q)$, $!A$ और $!B$ की भूमिकाएँ यहाँ क्रमशः
$\lnot \Box q \lif \Diamond \lnot p$, $\lnot\lnot p$ और~$p$ निभाते हैं।

जब किसी !!{formula} में उप-!!{formula} $\lnot\Diamond !A$ हो, तो हम
\olref{prop:rewriting} से उसे $\Box\lnot !A$ द्वारा बदल सकते हैं, क्योंकि
$\Log{K} \Proves \lnot\Diamond !A \liff \Box\lnot !A$। इसे और ऐसे ही
प्रतिस्थापनों को हम केवल ``$\lnot\Diamond$ के स्थान पर $\Box\lnot$'' से
दर्शाएँगे।

निम्न प्रतिज्ञप्ति यह उचित ठहराती है कि हम !!{derivability} के परिणाम
स्कीमागत रूप से स्थापित कर सकते हैं। उदाहरणार्थ, पूर्ववर्ती प्रतिज्ञप्ति
केवल $\Log{K} \Proves \lnot\Box p \lif \Diamond \lnot p$ नहीं, बल्कि
स्वेच्छ~$!A$ के लिए
$\Log{K} \Proves \lnot\Box !A \lif \Diamond\lnot !A$ स्थापित करती है।

\begin{prop}
  यदि $!A$, $!B$ का प्रतिस्थापन-दृष्टांत है और $\Log{K} \Proves !B$, तो
  $\Log{K} \Proves !A$।
\end{prop}

\begin{proof}
  यह जाँचना श्रमसाध्य पर नियमित है ($!B$ की !!{derivation} की लंबाई पर
  आगमन से) कि पूरी !!{derivation} पर प्रतिस्थापन लागू करने से भी सही
  !!{derivation} मिलती है। विशेषतः, सर्वसत्यतात्मक दृष्टांतों के
  प्रतिस्थापन-दृष्टांत स्वयं सर्वसत्यतात्मक दृष्टांत हैं;
  \iftag{prvDiamond}{\Dual{} और~}{}\Ax{K} के दृष्टांतों के
  प्रतिस्थापन-दृष्टांत स्वयं \iftag{prvDiamond}{\Dual{} और~}{}\Ax{K} के
  दृष्टांत हैं; और आधारवाक्य(यों) तथा निष्कर्ष दोनों में !!{propositional variable}ों
  के स्थान पर !!{formula} रखने पर \MP{} और~\Nec{} के अनुप्रयोग सही रहते हैं।
\end{proof}


\end{document}
